\subsection{Dataset}\label{sec:dataset}

El dataset es la carpeta \nolinkurl{data/dataset_bcs/}, que se sube al bucket de S3 del equipo:

\begin{verbatim}
dataset_bcs/
    imagenes/
        posterior/<id_vaca>/<foto>.jpg
        lateral_izq/<id_vaca>/<foto>.jpg
        lateral_der/<id_vaca>/<foto>.jpg
        descartadas/<id_vaca>/<foto>.jpg
    indice_imagenes.csv
    vacas.csv
    evaluaciones.csv
    tiempos.csv
    LEEME.md
\end{verbatim}

\begin{table}[H]
  \centering
  \caption{Resumen del dataset.}
  \label{tab:dataset}
  \begin{tabular}{ll}
    \toprule
    Concepto & Valor \\
    \midrule
    Fotos & 333 de 67 vacas \\
    Fotos para entrenar & 325 de 66 vacas \\
    Evaluaciones de BCS & 81 (79 de vacas que entran al entrenamiento) \\
    Pliegues & 5, separados por vaca \\
    Tamaño & 1.61 GB \\
    \bottomrule
  \end{tabular}
\end{table}

Las rutas de los CSV son relativas a \nolinkurl{dataset_bcs/}, así que al descargar el dataset hay que conservar esta estructura. El nombre del bucket no se escribe en el repositorio: los notebooks lo leen de la variable de entorno \texttt{S3\_BUCKET}. \textbf{Pendiente:} indicar en qué carpeta del bucket queda \nolinkurl{dataset_bcs/}.

\subsubsection*{Comparación entre pandas y PySpark}

El paso de integración, el diseño de prueba y la construcción de la etiqueta también los programamos con PySpark, en un contenedor Docker con Java 11 y Spark 3.5.0. Lo hicimos para practicar un enfoque que serviría si el volumen creciera. Los scripts de esa versión están en el bucket del equipo y no en el repositorio. \textbf{Pendiente:} decidir qué pasa con los scripts del bucket.

Las dos versiones generaron exactamente el mismo índice. Con este volumen, PySpark tardó unas 40 veces más en total, por el costo de arrancar Spark (mediana de 5 corridas, \nolinkurl{tiempos.csv}).

\begin{table}[H]
  \centering
  \caption{Tiempos de pandas y PySpark (mediana de 5 corridas).}
  \label{tab:tiempos}
  \begin{tabular}{lrrr}
    \toprule
    Versión & Total (s) & Arranque de sesión (s) & Procesamiento (s) \\
    \midrule
    pandas & 0.229 & No aplica & 0.027 \\
    PySpark (modo local) & 9.237 & 1.798 & 6.492 \\
    \bottomrule
  \end{tabular}
\end{table}

Por eso el notebook usa pandas y PySpark queda como comparación. \textbf{Pendiente:} definir si el procesamiento se pasa a PySpark en AWS, como pide el contexto del curso.
